\documentclass[10pt, aspectratio=169]{beamer}

% Theme selection
\usetheme{Madrid}
\usecolortheme{default}

% Packages
\usepackage[utf8]{inputenc}
\usepackage[T1]{fontenc}
\usepackage[vietnamese]{babel}
\usepackage{hyperref}
\usepackage{xcolor}

% Setup hyperref colors
\hypersetup{
    colorlinks=true,
    linkcolor=blue,
    urlcolor=blue,
}

% Author information extracted from TEMPLATE.tex
\title[Healthcare Data Lakehouse]{Onboarding Thông Tin Dự Án \& Quy Trình Làm Việc}
\subtitle{Healthcare Data Lakehouse: COVID-19 Analytics}
\author[Đỗ Kiến Hưng]{
    \textbf{Đỗ Kiến Hưng} \\
    Lead Data Engineer \\
    \footnotesize HCMUTE - Nhóm 01 \\
    \href{mailto:dkh1105.work@gmail.com}{dkh1105.work@gmail.com} $|$ \href{https://github.com/darktheDE}{github.com/darktheDE}
}
\institute[HCMUTE]{Đại học Sư phạm Kỹ thuật TP.HCM}
\date{\today}

\begin{document}

% Title slide
\begin{frame}
    \titlepage
\end{frame}

% TOC
\begin{frame}{Nội dung Onboarding}
    \tableofcontents
\end{frame}

% Section 1: Overview
\section{Tổng quan Dự Án}

\begin{frame}{Giới thiệu chung}
    \textbf{Mục tiêu:} 
    Xây dựng Data Lakehouse hoàn chỉnh, tách biệt Storage \& Compute, sử dụng 100\% công cụ Open-source, đảm bảo tính ACID và tự động hóa quy trình phân tích dữ liệu y tế (COVID-19).
    
    \vspace{0.3cm}
    \textbf{Phạm vi \& Dữ liệu:}
    \begin{itemize}
        \item Bộ dữ liệu mô phỏng bệnh án \textbf{Synthea™ COVID-19 100K} bệnh nhân.
        \item Khai thác 3 bảng chính: \texttt{patients}, \texttt{encounters}, \texttt{conditions}.
        \item Kết quả: Trực quan hóa Dashboard BI, API Data-as-a-Service cho ứng dụng nội bộ.
    \end{itemize}
\end{frame}

\begin{frame}{Kiến trúc \& Công nghệ (Docker Environment)}
    Dự án triển khai mô hình Medallion 5 lớp Data Lakehouse:
    \begin{itemize}
        \item \textbf{Data Source:} \texttt{PostgreSQL} (mô phỏng dữ liệu RDBMS bệnh viện).
        \item \textbf{Ingestion \& Orchestration:} \texttt{Apache Airflow} (điều phối ETL/ELT Batch).
        \item \textbf{Storage \& Table Format:} \texttt{MinIO} (Object Storage gốc) + \texttt{Apache Iceberg} (chuẩn bảng) + \texttt{Hive Metastore} (quản lý metadata).
        \item \textbf{Data Processing:} \texttt{Apache Spark} (xử lý dữ liệu Bronze $\rightarrow$ Silver $\rightarrow$ Gold).
        \item \textbf{Serving \& Consumption:} \texttt{Trino} (SQL Query tốc độ cao), \texttt{Apache Superset} (BI Dashboard) \& \texttt{FastAPI} (RESTful API).
    \end{itemize}
\end{frame}

% Section 2: Agile/Scrum Workflow
\section{Quy trình Công việc (Plane.so)}

\begin{frame}{Nhận \& Quản lý Task trên Plane.so}
    Mô hình làm việc \textbf{Agile/Scrum} ứng dụng \textbf{Plane.so}:
    \begin{enumerate}
        \item \textbf{Nhận Task:} Lựa chọn Task trên Plane.so, đọc \texttt{DEVELOPMENT.md} và yêu cầu chi tiết (Step-by-step required).
        \item \textbf{Phân công:} Tự Assign chính mình vào tác vụ (\textit{Assignee}) và cài đặt \textit{Reviewer} cho mình.
        \item \textbf{Workflow:} Chuyển trạng thái từ \texttt{Todo} sang \texttt{In Process} bắt đầu làm.
        \item \textbf{Deadline \& Họp:} 
        \begin{itemize}
            \item Phải cấu hình \textbf{Deadline} cho lập trình viên (thời điểm nộp PR).
            \item Meeting team định kỳ \textbf{2 ngày/buổi, lúc 22h30 mỗi tối}.
        \end{itemize}
    \end{enumerate}
\end{frame}

% Section 3: Git & Development
\section{Quy trình Lập trình \& PR}

\begin{frame}[fragile]{Quy trình Cấp phát Nhánh \& Coding}
    \textbf{1. Phân nhánh (Git Branching)} \\
    Tạo nhánh làm việc từ \texttt{develop} với cấu trúc chuẩn:
    \begin{verbatim}
    git checkout -b feature/<module>/<mã-task-trên-plane>
    \end{verbatim}
    
    \vspace{0.2cm}
    \textbf{2. Quản trị Bộ nhớ (Memory) \& Tài liệu} \\
    \textit{(Bắt buộc / Khuyến khích mạnh)}: Viết một file Markdown tên \texttt{<mã-task>.md} lưu lại:
    \begin{itemize}
        \item Các bước đã xây dựng / thiết lập.
        \item Lỗi \textit{(bugs)} gặp phải và cách khắc phục \textit{(fixes)}.
        \item Công dụng: Giảm thiểu lỗi cho AI Agent, làm tài liệu tham khảo chéo cho team và nguyên liệu viết Tutorial/Báo cáo cuối kỳ.
    \end{itemize}
\end{frame}

\begin{frame}{Quy trình Review \& Merge}
    \textbf{Sau khi phát triển xong, bắt đầu PR Workflow:}
    \begin{itemize}
        \item \textbf{Tạo PR:} Push nhánh feature và tạo \textit{Pull Request} vào \texttt{develop}.
        \item \textbf{Cài đặt Review Deadline:} Phân định số giờ tối đa cho Reviewer kiểm tra code và Accept PR.
        \item \textbf{Quy trình Duyệt 3 bên:}
        \begin{enumerate}
            \item Dựng môi trường test, Reviewer kiểm tra code (Comment chỉnh sửa nếu có).
            \item Reviewer \texttt{Approve}.
            \item Quản lý Repository kiểm tra chéo và thực hiện \texttt{Merge}.
        \end{enumerate}
    \end{itemize}
\end{frame}

\begin{frame}
    \begin{center}
        \Huge \textbf{Q\&A} \\
        \vspace{0.8cm}
        \normalsize Vui lòng tuân thủ chặt chẽ \texttt{DEVELOPMENT.md} \\
        \small Chúc mọi người có một học kỳ \& dự án thành công!
    \end{center}
\end{frame}

\end{document}
